\scanpage{23}
\uline{Lemme} 1 Soit $f:S'\to S$ un épimorphisme propre de schémas [\inserted{loc} noethériens], et $\deleted{\unsure{Z|X}}X$ un schéma \inserted{formellement} non ramifié sur $S$.
\margin{[ ] inutile si $X$ \uline{loc} de t.f.?}
Alors la suite
\[
\operatorname{Hom}_S(S,X)\to\operatorname{Hom}_S(S',X)\rightrightarrows\operatorname{Hom}_{S}(S'',X)
\]
est exacte, où $S''=S'\times_S S'$.

NB. Le fait que $f$ est un épimorphisme s'exprime par la condition que l'hom
\[
\underline O_S\to f_\ast\mathcal O_{S'}
\]
est \uline{injectif}.

Soit
\[
g':S'\to X
\]
tel que $g'\operatorname{pr}_1=g'\operatorname{pr}_2$, à montrer que $g'$ se factorise par $S$. Tout d'abord il se factorise par $S$ du point de vue application continue, soit $g_0$, et $f^{-1}(g_0^{-1}(\underline O_X))\simeq {g'}^{-1}(O_X)\to\mathcal O_{S'}$ définit un hom
\[
g_0^{-1}(\underline O_X)\to f_\ast(\underline{\mathcal O}_{S'})
\]
;
\diagram{\begin{tikzcd}[ampersand replacement=\&,column sep=small,row sep=small]
\&\& X \arrow[d,shift right]\\
S'' \arrow[r,shift left,"p_1"] \arrow[r,shift right,"p_2"'] \arrow[rr,bend right=25,"h"'] \& S' \arrow[ur,"g'"] \arrow[r,"f"'] \& S \arrow[u,dotted,shift right,"g_0"']
\end{tikzcd}}
l'hyp. $g'p_1=g'p_2$ s'exprime par le fait que $g_0^{-1}(\underline O_X)$ s'envoie dans
\[
\mathcal A=\operatorname{Ker}\bigl(f_\ast(\underline{\mathcal O}_{S'})\rightrightarrows\deleted{\unsure{g_\ast}}h_\ast(\underline O_{S''})\bigr)
\]
donc $g'$ se factorise par
\[
S'\to\operatorname{Spec}(\mathcal A)=S'_1\to X.
\]

\margin{\uline{lemme}}
Or EGA III 3 indique aisément que $S'_1\to S'$ est \inserted{fini +} \uline{radiciel}, c'est d'ailleurs un épimorphisme car $\underline O_S\hookrightarrow\deleted{\unsure{f_\ast\mathcal O_{S'}|f_\ast(\mathcal O_{S'})}}\mathcal A\hookrightarrow f_\ast(\underline{\mathcal O}_{S'})$, $f$ étant un épimorphisme. On est donc ramené à prouver le

\uline{Corollaire 1} Soit $f:S'\to S$ un épimorphisme de schémas \deleted{\unsure{noethériens}}; $f$ \deleted{fini} \inserted{entier} radiciel et $X$ \deleted{non ramifié} \inserted{\deleted{loc de t.f.} form\textsuperscript{t} net} sur $S$; \deleted{\unsure{Alors si $f$ fini}} \inserted{on suppose $f$ fini et $S$ loc. noeth, ou $X\to S$ loc. de t.f.} Alors
\[
\operatorname{Hom}_S(S,X)\to\operatorname{Hom}_S(S',X)\quad\text{est \uline{bijectif}.}
\]
% archive sheet number: 22

